\documentclass[14pt,a4paper]{extarticle}

\usepackage{fontspec}
\usepackage{polyglossia}
\usepackage{geometry}
\usepackage{setspace}
\usepackage{hyperref}
\usepackage{xurl}
\usepackage{enumitem}
\usepackage{titlesec}
\usepackage{microtype}
\usepackage{graphicx}
\usepackage{float}
\usepackage{caption}

\setmainlanguage{russian}
\setotherlanguage{english}
\setmainfont{Times New Roman}
\setmonofont{Consolas}
\newfontfamily\cyrillicfont{Times New Roman}
\newfontfamily\cyrillicfonttt{Consolas}

\geometry{
  left=30mm,
  right=15mm,
  top=20mm,
  bottom=20mm
}

\onehalfspacing
\setlist{nosep}
\setlength{\parindent}{1.25cm}
\setlength{\parskip}{0pt}
\emergencystretch=1.5em
\Urlmuskip=0mu plus 2mu

\titleformat{\section}
  {\normalfont\bfseries\large}
  {\thesection.}
  {0.5em}
  {}

\titleformat{\subsection}
  {\normalfont\bfseries}
  {\thesubsection.}
  {0.5em}
  {}

\hypersetup{
  colorlinks=false,
  pdfborder={0 0 0}
}

\newcommand{\screenshotfigure}[3]{
  \begin{figure}[H]
    \centering
    \IfFileExists{#1}{
      \includegraphics[width=#2\textwidth]{#1}
    }{
      \fbox{
        \begin{minipage}[c][0.28\textheight][c]{\dimexpr#2\textwidth-2\fboxsep-2\fboxrule\relax}
          \centering
          \textbf{Место для скриншота}\\[0.4cm]
          \texttt{#1}
        \end{minipage}
      }
    }
    \caption{#3}
  \end{figure}
}

\begin{document}

\pagenumbering{gobble}

\begin{titlepage}
  \begin{center}
    \vspace*{2cm}

    {\large ОТЧЕТ ПО ПРОЕКТУ\par}

    \vspace{2cm}

    {\Large\bfseries Cold War History: интерактивный проект о шифрах Холодной войны\par}

    \vspace{2.5cm}

    \begin{flushleft}
      \textbf{Участники проекта:}

      Неганов Михаил Юрьевич, группа: 1.2

      Данилов Максим Юрьевич, группа: 1.2

      Ишкаев Александр Олегович, группа: 1.2

      \vspace{0.5cm}

      \textbf{Тимлид проекта:} Неганов Михаил Юрьевич
      \end{flushleft}

    \vfill

    {\large 2026}

    \vspace{0.8cm}

    \small Репозиторий проекта: \url{https://github.com/ColdWar-History/History}
  \end{center}
\end{titlepage}

\pagenumbering{arabic}

\section{Синопсис проекта}

«Cold War History: интерактивный проект о шифрах Холодной войны» представляет собой образовательный веб-проект, объединяющий исторический материал, каталог классических шифров, криптографическую лабораторию и игровые задания. Пользователь может изучать события Холодной войны, знакомиться с историческими способами шифрования, шифровать и расшифровывать сообщения, выполнять тренировочные задания, проходить ежедневные вызовы и игровой режим «Инспектор связи».

Проект разработан как MVP, включающий frontend-приложение, backend на C\#/.NET, базу данных PostgreSQL с миграциями и seed-данными, а также Docker Compose для локального запуска. Основная идея проекта заключается в том, чтобы показать связь между историей, разведкой, коммуникациями и криптографией через интерактивный пользовательский опыт.

\section{Целевая аудитория и актуальность проекта}

Проект ориентирован прежде всего на школьников старших классов и студентов, интересующихся историей, обществознанием, международными отношениями, информатикой и базовой криптографией. Для этой аудитории особенно важен формат, в котором исторический материал не только читается, но и применяется на практике.

Актуальность проекта состоит в том, что Холодная война является одним из ключевых периодов истории XX века, а тема защищенной коммуникации, разведки и анализа информации остается понятной и современной для цифровой аудитории. Интерактивный формат помогает воспринимать исторические события не как набор дат, а как систему политических решений, информационного противоборства и технологических инструментов.

Проект может быть полезен:

\begin{itemize}
  \item школьникам и студентам как наглядный образовательный ресурс по истории и криптографии;
  \item преподавателям как демонстрационный материал для занятий по истории, информатике и проектной деятельности;
  \item пользователям, интересующимся шифрами, как безопасная среда для знакомства с учебными историческими алгоритмами;
  \item проектным командам как пример связки frontend, backend, базы данных и содержательного исторического материала в одном MVP.
\end{itemize}

\section{Описание содержательной исторической части}

Содержательная часть проекта построена вокруг связи между историческими событиями Холодной войны, коммуникациями, разведкой и средствами защиты информации. В проект включены ключевые события периода 1948--1989 годов: Берлинский кризис 1948 года, создание НАТО, Корейская война, смерть Сталина, создание Варшавского договора, Венгерское восстание, запуск «Спутника-1», инцидент с U-2, строительство Берлинской стены, Карибский кризис, операция «Анадырь», Пражская весна, разрядка и ОСВ-1, Афганская война, падение Берлинской стены.

Для каждого события предусмотрены дата, регион, тема, краткое описание, участники и связь с коммуникацией, разведкой или контролем информации. Такой подход позволяет изучать события не изолированно, а через общую логику политического противостояния и информационной безопасности.

Криптографическая часть проекта включает пять учебных алгоритмов: шифр Цезаря, Атбаш, шифр Виженера, Rail Fence и столбцовую перестановку. Для каждого алгоритма указаны описание, пример использования, параметры, пошаговое объяснение преобразования и ограничения применения. Алгоритмы представлены как исторические и учебные, а не как современные средства защиты данных.

\section{Процесс работы над историческим материалом}

Работа над исторической частью выполнялась по следующему алгоритму:

\begin{enumerate}
  \item Был определен временной диапазон проекта: от начала послевоенного противостояния до завершения Холодной войны.
  \item Команда выделила ключевые события, которые хорошо отражают политические кризисы, военные блоки, разведывательную деятельность и информационное противостояние.
  \item Для каждого события были подготовлены краткие описания, понятные образовательной аудитории.
  \item Исторические события были связаны с темами коммуникации, разведки, шифрования или контроля информации.
  \item Были выбраны простые учебные шифры, которые можно реализовать в MVP и объяснить без специальной математической подготовки.
  \item Материалы были распределены по карточкам, таймлайну и тематическим подборкам.
\end{enumerate}

При подготовке материалов использовались энциклопедические, архивные и музейные источники: Encyclopaedia Britannica, Wilson Center Digital Archive, материалы Wilson Center Cold War Project, публикации NSA/CSS Center for Cryptologic History, материалы National Cryptologic Museum и CIA FOIA Reading Room.

\section{Описание технической части проекта}

Проект реализован как веб-приложение с разделением на frontend, backend и database. Такой подход позволил разделить пользовательский интерфейс, бизнес-логику и хранение данных.

\subsection{Frontend}

Frontend реализован как SPA-приложение на React, TypeScript и Vite. Интерфейс включает главную страницу, каталог шифров, страницы отдельных шифров, исторический таймлайн, криптолабораторию, тренировочный режим, ежедневный вызов, игровой режим «Инспектор связи», профиль пользователя, лидерборд, вход в систему и редакторский раздел.

Клиентская часть обращается к backend через единый API gateway. В интерфейсе предусмотрены формы для работы с шифрами, отображение результата преобразования, шагов алгоритма и ограничений шифров на русском языке. В процессе доработки frontend был очищен от технических dev-надписей, чтобы пользовательский интерфейс выглядел как законченный продукт.

\subsection{Backend}

Backend написан на C\#/.NET и организован с учетом принципов гексагональной архитектуры. Для сервисов выделены слои API, Application, Domain и Infrastructure. Такое разделение помогает отделить бизнес-логику от инфраструктурных деталей и упростить тестирование.

В backend реализованы следующие сервисы:

\begin{itemize}
  \item Gateway: единая точка входа для frontend и маршрутизация запросов к сервисам;
  \item Auth: регистрация, вход, refresh-токены, роли пользователя, редактора и администратора;
  \item Crypto: шифрование, дешифрование, валидация параметров, пошаговое объяснение алгоритмов и ограничения шифров;
  \item Content: карточки шифров, исторические события, тематические подборки и версии материалов;
  \item Game: тренировочные задания, ежедневный вызов и режим «Инспектор связи»;
  \item Progress: история операций, очки, достижения, пользовательские метрики и лидерборд.
\end{itemize}

Для backend подготовлены тестовые harness-проекты, которые проверяют криптографические алгоритмы и ключевые application-сценарии auth, content, game и progress. Также подготовлен Docker Compose для локального запуска backend-сервисов, мигратора и PostgreSQL.

\subsection{Database}

База данных построена на PostgreSQL. Используется подход schema-per-service: для разных зон ответственности выделены схемы auth, content, game и progress. Такой подход соответствует микросервисной структуре backend и упрощает сопровождение данных.

В database-части подготовлены миграции up/down, seed-данные, ERD, а также операционные SQL-скрипты для backup, restore, мониторинга и retention-политики истории прогресса. Seed-данные включают администратора, исторические события, карточки шифров и тематические подборки.

\section{Оценка результатов}

В результате работы был подготовлен рабочий MVP образовательного проекта. Пользователь может изучать исторические события, открывать карточки шифров, выполнять преобразования в криптолаборатории, проходить игровые задания и сохранять прогресс после авторизации. Команда также подготовила техническую основу для дальнейшего расширения проекта: backend-сервисы, PostgreSQL-схему, миграции, тесты, Docker Compose и frontend-интерфейс.

К сильным сторонам проекта можно отнести связь исторического содержания с практическими заданиями, разделение архитектуры на понятные зоны ответственности, наличие seed-данных, тестов и локального сценария запуска. Проект не ограничивается справочным материалом, а предлагает пользователю взаимодействовать с темой через шифры и игровые сценарии.

Главное ограничение текущей версии заключается в отсутствии полноценного публичного хостинга всего проекта. Статический frontend можно разместить на GitHub Pages, однако backend и база данных пока рассчитаны на локальный запуск через Docker Compose. Также в MVP представлен ограниченный набор исторических материалов и учебных шифров. В дальнейшем проект можно развивать через облачный деплой backend/database, расширение исторической базы, добавление новых алгоритмов, улучшение игровых сценариев и полноценные end-to-end тесты.

Ссылка на репозиторий проекта: \url{https://github.com/ColdWar-History/History}

\section{Вклад участников}

Оценка вклада указана как процент выполнения каждым участником своей зоны ответственности.

\textbf{Неганов Михаил Юрьевич: 100\%.}
Отвечал за backend-часть и участвовал в подготовке исторического содержания. Реализовывал архитектуру backend на C\#/.NET, сервисы gateway, auth, crypto, content, game и progress, криптографические алгоритмы, тесты, Docker Compose и связку backend с историческим материалом.

\textbf{Данилов Максим Юрьевич: 100\%.}
Отвечал за database-часть и техническое задание. Прорабатывал требования проекта, структуру базы данных, схемы auth/content/game/progress, миграции, seed-данные, ERD, а также SQL-скрипты для backup, restore и мониторинга.

\textbf{Ишкаев Александр Олегович: 100\%.}
Отвечал за frontend-часть и дизайн сайта. Реализовывал React/TypeScript/Vite-приложение, пользовательские страницы, криптолабораторию, игровые режимы, профиль, авторизацию, редакторский раздел, внешний вид интерфейса и пользовательский сценарий.

\section{Скриншоты проекта}

Ниже представлены основные экраны пользовательского интерфейса. Скриншоты демонстрируют главную страницу, навигацию по историческим материалам, каталог шифров, криптолабораторию и ежедневный игровой вызов.

\screenshotfigure{screenshots/home.png}{0.98}{Главная страница проекта с навигацией, кратким описанием продукта и счетчиками опубликованных материалов.}

\screenshotfigure{screenshots/catalog.png}{0.98}{Каталог шифров с фильтрами по поиску, категории и эпохе. Пользователь может перейти к досье шифра или сразу открыть выбранный алгоритм в криптолаборатории.}

\screenshotfigure{screenshots/timeline.png}{0.98}{Исторический таймлайн с фильтрацией по региону, году и теме, а также тематическими коллекциями событий.}

\screenshotfigure{screenshots/crypto-lab.png}{0.98}{Криптолаборатория: выбор алгоритма, ввод сообщения, параметры шифрования, результат и пошаговое объяснение преобразования.}

\screenshotfigure{screenshots/daily.png}{0.98}{Ежедневный вызов: короткое задание по одному из шифров и форма для отправки решения.}

\section{Источники и материалы}

\begin{enumerate}
  \item Wilson Center Digital Archive. Declassified documents and educational materials on Cold War history. \url{https://www.wilsoncenter.org/digital-archive}
  \item Wilson Center. About the Cold War Project. \url{https://www.wilsoncenter.org/about-18}
  \item NSA/CSS Center for Cryptologic History. Historical Publications. \url{https://www.nsa.gov/History/Cryptologic-History/Historical-Publications/}
  \item National Cryptologic Museum. About the Museum. \url{https://www.nsa.gov/History/National-Cryptologic-Museum/About-the-Museum/}
\end{enumerate}

\end{document}
